% \usepackage{showkeys}

%%%%%%%%%%%%%%%%%%%%%%%%%
% Line numbering
\usepackage[left,mathlines,pagewise]{lineno}
\newcommand*\patchAmsMathEnvironmentForLineno[1]{%
\expandafter\let\csname old#1\expandafter\endcsname\csname #1\endcsname
\expandafter\let\csname oldend#1\expandafter\endcsname\csname end#1\endcsname
\renewenvironment{#1}%
{\linenomath\csname old#1\endcsname}%
{\csname oldend#1\endcsname\endlinenomath}}%
\newcommand*\patchBothAmsMathEnvironmentsForLineno[1]{%
\patchAmsMathEnvironmentForLineno{#1}%
\patchAmsMathEnvironmentForLineno{#1*}}%
\AtBeginDocument{%
\patchBothAmsMathEnvironmentsForLineno{equation}%
\patchBothAmsMathEnvironmentsForLineno{align}%
\patchBothAmsMathEnvironmentsForLineno{flalign}%
\patchBothAmsMathEnvironmentsForLineno{alignat}%
\patchBothAmsMathEnvironmentsForLineno{gather}%
\patchBothAmsMathEnvironmentsForLineno{multline}%
}


%very horrible equation
\usepackage{array} %\displaystyle%
\newenvironment{ma}{\renewcommand{\arraystretch}{2.0}\begin{array}{>{\displaystyle}r >{\displaystyle}c >{\displaystyle}l}}{\end{array}\renewcommand{\arraystretch}{1.0}}%

\usepackage{amsmath,amsxtra,amsthm}
\usepackage[varg]{txfonts}
\usepackage[unicode]{hyperref}
\usepackage{framed}
%\usepackage[amsmath,hyperref,framed,thmmarks]{ntheorem}
\usepackage{mathrsfs}
\usepackage{graphicx}

%\usepackage{bbm}

%%%%%%%%%%%%%%%%%%%%%%%%%
 %vruler
%\usepackage{vruler}
%\setvruler[10pt][0][2][4][1][1cm][1cm][-5mm][\textheight]

%%%%%%%%%%%%%%%%%%%%%%%%%
% Theorems

\makeatletter
\def\th@plain{%
%%  \let\thm@indent\noindent % no indent
\thm@headfont{\bfseries}% heading font is bold
\thm@notefont{}% same as heading font
%%  \thm@headpunct{.}% add period after heading
%%  \let\thm@swap\@gobble
%%  \thm@preskip\topsep
%%  \thm@postskip\theorempreskipamount
  \itshape % body font
}
\makeatother

\newtheorem{theorem}{Theorem}%
\renewcommand{\thetheorem}{\Roman{theorem}}
\newtheorem{lemma}{Lemma}[section]%
\newtheorem{proposition}[lemma]{Proposition}%
\newtheorem{corollary}[lemma]{Corollary}%

\newcounter{mythmct}
\setcounter{mythmct}{0}
%\renewcommand{\themythmct}{\Roman{mythmct}}
\newtheorem{mythm}[mythmct]{Theorem}

\newtheorem{remark}[lemma]{Remark}
\newtheorem{example}[lemma]{Example}%


%%%%%%%%%%%%%%%%%%%%%%%%%
% Hyperlinks in blue

\hypersetup{
breaklinks=true,
colorlinks=true,
linkcolor=blue,
citecolor=blue,
urlcolor=blue,
}

%%%%%%%%%%%%%%%%%%%%%%%%%
% Paragraph layout
\setlength{\parindent}{0ex}
\setlength{\parskip}{1ex}

%%%%%%%%%%%%%%%%%%%%%%%%%
% displaybreaks
\allowdisplaybreaks

%%%%%%%%%%%%%%%%%%%%%%%%%
% boldmath (from http://www.volkerschatz.com/tex/thacks.html)
\def\bfseries{\fontseries \bfdefault \selectfont \boldmath}

%%%%%%%%%%%%%%%%%%%%%%%%%
% Roman numbers for enumerate
\renewcommand{\theenumi}{\textup{(\roman{enumi})}}
\renewcommand{\labelenumi}{\theenumi}

%%%%%%%%%%%%%%%%%%%%%%%%%
% fnsymbols by asterisks
\makeatletter
\def\@fnsymbol#1{\ensuremath{\ifcase#1\or *\or **\or {**}* \or {**}{**}\else\@ctrerr\fi}}
\makeatother

\setcounter{tocdepth}{1}